\section{B\'usqueda binaria}
\whenuse{Existe una condici\'on mon\'otona falso--verdadero o verdadero--falso y
se busca el primer/\'ultimo valor que la satisface, incluso sobre la respuesta.}

Busca la primera posici\'on donde una condici\'on mon\'otona pasa de falsa a
verdadera. 


\complexity{Dominio discreto de $n$ candidatos: $O(\log(n+1))$ evaluaciones;
si cada chequeo cuesta $C$, tiempo $O(C\log(n+1))$. Para reales con ancho
inicial $W$ y tolerancia $\varepsilon$, $O(\log(1+W/\varepsilon))$
evaluaciones hasta alcanzar la tolerancia o el l\'imite de precisi\'on.
Memoria adicional $O(1)$ aparte del chequeo.}

\begin{lstlisting}[style=cpp]
int l = start;
int r = end;
while (l <= r) {
    int m = l + (r - l) / 2; 
    if (can(m)) { 
        // La condicion se cumple en m, ajustar el limite superior 
        r = m - 1; 
    } else { 
        // La condicion no se cumple en m, ajustar el limite inferior 
       l = m + 1; 
    } 
}
// Esto es en caso de que la condicion inicialmente sea falsa y luego verdadera. Si la condicion es verdadero--falso, se debe invertir los contenidos del if-else
\end{lstlisting}

El algoritmo anterior permite usar la busqueda binaria unicamente en un dominio discreto.
En este caso $r$ nos da el ultimo valor que no cumple la condicion y $l$ el menor valor que si la cumple.
Para reales, mantener un intervalo con \texttt{can(l)==false} y
\texttt{can(r)==true}. Usar la tolerancia pedida por el enunciado y cortar
si el punto medio coincide con un extremo para evitar ciclos por redondeo.

\begin{lstlisting}[style=cpp]
double l = start;
double r = end;
double eps = 1e-6; // Ajustar al error permitido por el enunciado.
while (r - l > eps) {
    double m = l + (r - l) / 2.0;
    if (m == l || m == r) break; // Limite de precision de double.
    if (can(m)) { 
        // La condicion se cumple para m, ajustamos el limite superior 
        r = m; 
    } else { 
       // La condicion no se cumple para m, ajustamos el limite inferior 
       l = m; 
    } 
}
// r cumple la condicion; l no la cumple.
\end{lstlisting}

El ancho se divide por dos en cada iteraci\'on. La precisi\'on final tambi\'en
est\'a limitada por la representaci\'on de \texttt{double}: el corte por
redondeo puede detener la b\'usqueda antes de alcanzar \texttt{eps}.
Si se permite error relativo, adaptar la tolerancia a la escala del resultado.
